% ==========================================================================
%  Fragebogen / Interview-Leitfaden – Vorlage
%  Build:  latexmk LeitfadenFragekatalog.tex   (im Ordner .latex)
% ==========================================================================
\documentclass[11pt,a4paper]{article}
\usepackage[T1]{fontenc}
\usepackage[utf8]{inputenc}
\usepackage[ngerman]{babel}
\usepackage{lmodern}
\usepackage[margin=2.2cm]{geometry}
\usepackage{amssymb}        % \square
\usepackage{tabularx}
\usepackage{xcolor}
\usepackage{fancyhdr}

\setlength{\parindent}{0pt}
\setlength{\parskip}{4pt}

% --------------------------------------------------------------------------
%  Anpassen
% --------------------------------------------------------------------------
\newcommand{\projekt}{Ground Speed Sensor}
\newcommand{\dokumenttitel}{Leitfaden -- Kundeninterview}
\newcommand{\version}{v0.1}

% --------------------------------------------------------------------------
%  Kopf-/Fusszeile
% --------------------------------------------------------------------------
\pagestyle{fancy}
\fancyhf{}
\lhead{\small\projekt}
\rhead{\small\dokumenttitel}
\lfoot{\small\version}
\rfoot{\small Seite \thepage}
\renewcommand{\headrulewidth}{0.4pt}

\definecolor{akzent}{RGB}{0,84,140}
% \platzbedarf{Höhe}: Seitenumbruch, falls weniger als <Höhe> Platz bleibt
\newcommand{\platzbedarf}[1]{\par\vskip 0pt plus #1\penalty-200\vskip 0pt plus -#1\relax}
\makeatletter
\renewcommand{\section}{\platzbedarf{7cm}\kern0pt\@startsection{section}{1}{0pt}%
  {14pt plus 2pt}{6pt}{\large\bfseries\color{akzent}}}
\makeatother

% --------------------------------------------------------------------------
%  Befehle für den Fragebogen
% --------------------------------------------------------------------------
\newcounter{frage}

% \frage{Text}                 – nummerierte Frage
\makeatletter
\newcommand{\frage}[1]{%
  \par\if@nobreak\else\platzbedarf{5cm}\fi % Frage nicht vom Antwortfeld trennen
  \medskip\refstepcounter{frage}%
  \noindent\textbf{F\thefrage.}\hspace{0.5em}#1\par\nopagebreak}
\makeatother

% \hinweis{Text}               – Nachfrage / Hinweis für Interviewer (grau, kursiv)
\newcommand{\hinweis}[1]{{\small\itshape\color{gray} $\rightarrow$ #1}\par}

% \linien{n}                   – n Antwortlinien
\makeatletter
\newcommand{\linien}[1]{%
  \par\nopagebreak
  \@tempcnta=#1\relax
  \loop\ifnum\@tempcnta>0
    \vspace{0.75cm}\noindent\rule{\linewidth}{0.3pt}\par
    \advance\@tempcnta by -1
  \repeat}
\makeatother

% \auswahl{Ja \item Nein \item ...}  – Ankreuzoptionen nebeneinander
\newcommand{\kaestchen}{$\square$~}
\newcommand{\auswahl}[1]{%
  \par\nopagebreak\hspace*{1em}%
  {\def\item{\hspace{1.5em}\kaestchen}\kaestchen #1}\par}

% \skala{links}{rechts}        – Bewertungsskala 1–5
\newcommand{\skala}[2]{%
  \par\nopagebreak\smallskip
  \begin{tabularx}{\linewidth}{@{}X*{5}{c}X@{}}
    & 1 & 2 & 3 & 4 & 5 & \\
    \raggedleft\small #1 & \kaestchen & \kaestchen & \kaestchen
      & \kaestchen & \kaestchen & \small #2
  \end{tabularx}\par}

% ==========================================================================
\begin{document}

{\LARGE\bfseries\color{akzent} \dokumenttitel}\par
\medskip
{\large \projekt}\par
\bigskip

% ---------------------------- Interviewdaten ------------------------------
\renewcommand{\arraystretch}{1.6}
\begin{tabularx}{\linewidth}{@{}l X l X@{}}
  Datum:        & \hrulefill & Interviewer/in: & \hrulefill \\
  Unternehmen:  & \hrulefill & Gesprächspartner/in: & \hrulefill \\
  Funktion:     & \hrulefill & Dauer:          & \hrulefill \\
\end{tabularx}
\renewcommand{\arraystretch}{1}

\medskip
\fbox{\parbox{\dimexpr\linewidth-2\fboxsep-2\fboxrule}{\small
  \textbf{Einleitung (vorlesen):} Vielen Dank, dass Sie sich Zeit nehmen.
  Ziel des Gesprächs ist es, Ihre Anforderungen und Erfahrungen zu verstehen.
  Es gibt keine richtigen oder falschen Antworten. Die Angaben werden
  vertraulich behandelt.}}

% ==========================================================================
\section{Unternehmen und Anwendung}

\frage{Beschreiben Sie kurz Ihr Unternehmen und Ihre Rolle.}
\linien{2}

\frage{In welchen Anwendungen messen Sie heute Geschwindigkeit über Grund?}
\hinweis{Fahrzeugtyp, Einsatzumgebung, Stückzahlen}
\linien{3}

\frage{Welche Branche trifft am ehesten zu?}
\auswahl{Landwirtschaft \item Bau \item Bahn \item Automotive \item Andere: \rule{3cm}{0.3pt}}

% ==========================================================================
\section{Heutige Lösung}

\frage{Welche Sensoren oder Verfahren setzen Sie aktuell ein?}
\auswahl{Radar \item GNSS \item Rad-Encoder \item Optisch \item Keine}
\linien{1}

\frage{Wie zufrieden sind Sie mit der heutigen Lösung?}
\skala{sehr unzufrieden}{sehr zufrieden}

\frage{Was sind die grössten Probleme der heutigen Lösung?}
\hinweis{Genauigkeit, Kosten, Robustheit, Integration, Wartung?}
\linien{3}

% ==========================================================================
\section{Anforderungen}

\frage{Wie wichtig sind Ihnen folgende Eigenschaften?}
\smallskip
\renewcommand{\arraystretch}{1.3}
\begin{tabularx}{\linewidth}{@{}X*{5}{>{\centering\arraybackslash}p{0.9cm}}@{}}
  \small\itshape 1 = unwichtig, 5 = sehr wichtig & 1 & 2 & 3 & 4 & 5 \\
  Genauigkeit           & \kaestchen & \kaestchen & \kaestchen & \kaestchen & \kaestchen \\
  Preis                 & \kaestchen & \kaestchen & \kaestchen & \kaestchen & \kaestchen \\
  Robustheit (Staub, Wasser, Vibration) & \kaestchen & \kaestchen & \kaestchen & \kaestchen & \kaestchen \\
  Baugrösse / Gewicht   & \kaestchen & \kaestchen & \kaestchen & \kaestchen & \kaestchen \\
  Einfache Integration  & \kaestchen & \kaestchen & \kaestchen & \kaestchen & \kaestchen \\
\end{tabularx}
\renewcommand{\arraystretch}{1}

\frage{Welcher Geschwindigkeitsbereich und welche Genauigkeit werden benötigt?}
\linien{2}

\frage{Welche Schnittstellen sind erforderlich?}
\auswahl{CAN \item RS-232/485 \item Analog \item Puls/Frequenz \item Ethernet}

% ==========================================================================
\section{Wirtschaftliches}

\frage{Welchen Preis pro Sensor halten Sie für angemessen?}
\linien{1}

\frage{Wer entscheidet bei Ihnen über die Beschaffung?}
\linien{2}

% ==========================================================================
\section{Abschluss}

\frage{Gibt es etwas, das wir nicht angesprochen haben, das Ihnen wichtig ist?}
\linien{3}

\frage{Dürfen wir Sie für Rückfragen oder einen Prototyp-Test erneut kontaktieren?}
\auswahl{Ja \item Nein}

\vfill
\small\textit{Notizen des Interviewers:}
\linien{3}

\end{document}
